% !TeX root = ../main.tex
\section{Постановка задачи и исходные данные}
В соответствии с вариантом №\,8 на вход цифрового приемника поступает дискретная выборка $Y_1^N = \{y_k\}$, $k = 0, \dots, N-1$ наблюдений сигнала ГЛОНАСС L1OF, загруженная из файла \texttt{Yvar8.txt}. Длительность реализации составляет $T = 1\text{ мс}$ на промежуточной частоте $f_0 = 8\text{ МГц}$ при частоте дискретизации $F_d = 33\text{ МГц}$. Число отсчетов в реализации:
\begin{equation}
N = T \cdot F_d = 10^{-3} \cdot 33 \cdot 10^6 = 33000.
\end{equation}
Модель сигнала на выходе 3-уровневого АЦП ($y_k \in \{-1, 0, +1\}$):
\begin{equation}
y_k = a \cdot G(k T_d - \tau_{\text{з}}) \cos\bigl(2\pi(f_0 + f_{\text{д}}) k T_d + \varphi_0\bigr) + n_k,
\end{equation}
где:
\begin{itemize}
    \item $n_k$ --- дискретный белый гауссовский шум (ДБГШ) с единичной дисперсией $\sigma_n^2 = 1$;
    \item $a$ --- неизвестная постоянная амплитуда, распределенная по закону Рэлея с параметром $\sigma_a = 0.06$ ($q_{c/n_0} = 45\text{ дБ}\cdot\text{Гц}$);
    \item $\varphi_0 \in [0, 2\pi)$ --- неизвестная начальная фаза с равномерным распределением;
    \item $\tau_{\text{з}} \in [0, T_{\text{дк}})$ --- неизвестная задержка сигнала ($T_{\text{дк}} = 1\text{ мс}$);
    \item $f_{\text{д}} \in [-2000, +2000]\text{ Гц}$ --- неизвестный доплеровский сдвиг частоты;
    \item $G(t)$ --- периодическая псевдослучайная последовательность (ПСП) дальномерного кода ГЛОНАСС длительностью 511 чипов со скоростью $511\text{ кбит/с}$ (длительность чипа $\tau_c = 1/511\text{ мс}$).
\end{itemize}

Область поиска дискретизируется на сетку из $1022 \times 9 = 9198$ ячеек:
\begin{itemize}
    \item по задержке $\tau_{\text{з}}$: 1022 шага в диапазоне $[0, 1)\text{ мс}$ с шагом $\Delta \tau = 1/1022\text{ мс} \approx 0.9785\text{ мкс}$ (по 2 отсчета на один чип ПСП);
    \item по частоте $f_{\text{д}}$: 9 значений от $-2000\text{ Гц}$ до $+2000\text{ Гц}$ с шагом $\Delta f = 500\text{ Гц}$.
\end{itemize}
